\chapter{Geographic Tongue}

\section*{Definition}
Geographic tongue (benign migratory glossitis) is a benign inflammatory condition of the tongue characterized by erythematous, depapillated patches surrounded by slightly elevated whitish borders that change in size, shape, and location over days to weeks.

\section*{What Happens in Your Body}
The underlying mechanism involves localized desquamation of the filiform papillae from the dorsal tongue surface, likely mediated by T-lymphocyte infiltration and immune signaling proteins release at the epithelial level. The migratory nature reflects cycles of inflammation, healing, and recurrence in different areas.

\section*{Causes}
\begin{itemize}
  \item \textbf{Idiopathic:} No identifiable cause in the majority of cases
  \item \textbf{Genetic predisposition:} HLA-DR5, HLA-Cw6 association; familial clustering suggests autosomal dominant inheritance with incomplete penetrance
  \item \textbf{Associated conditions:} Psoriasis (10\% prevalence), fissured tongue (co-occurrence), atopic conditions (allergic rhinitis, asthma, eczema)
  \item \textbf{Potential triggers:} Emotional stress, spicy or acidic foods, hormonal fluctuations (menstruation), vitamin B deficiency
\end{itemize}

\section*{When to See a Doctor}
See your doctor if:
\begin{itemize}
  \item Lesions cause significant burning or pain, especially with certain foods
  \item New or changing oral lesions in you > 40 years old (exclude malignancy)
  \item Suspected oral psoriasis or concurrent skin lesions
  \item Ulceration, bleeding, or induration (not typical of geographic tongue)
  \item Associated systemic symptoms (fever, weight loss, lymphadenopathy)
\end{itemize}

\section*{Related Symptoms}
\begin{itemize}
  \item Migratory erythematous patches on the tongue dorsum and lateral borders
  \item Burning or sensitivity to spicy, salty, or acidic foods
  \item Fissured tongue (scrotal tongue) in up to 50\% of people
  \item Usually asymptomatic; incidental finding on oral examination
  \item Rarely, involvement of other oral mucosal sites (ectopic geographic tongue)
\end{itemize}
\textbf{Important:} Geographic tongue affects 1--3\% of the population, is entirely benign, and requires no treatment in asymptomatic people; the key differential includes oral candidiasis, lichen planus, and erythematous candidiasis.

\section*{Self-Care / Home Management}
\begin{itemize}
  \item Avoidance of trigger foods: spicy, acidic (citrus, tomatoes), hot, or salty foods
  \item Avoid alcohol-based mouthwashes and astringent toothpastes
  \item Topical anesthetic mouth rinses (viscous lidocaine or diphenhydramine elixir) before meals if painful
  \item Sucralfate suspension as a mouth rinse to create a protective barrier
  \item Stress reduction techniques (yoga, meditation) if stress is a trigger
  \item Zinc or vitamin B complex supplementation if dietary deficiency suspected
\end{itemize}

\section*{How Your Doctor Will Evaluate You}
\begin{itemize}
  \item Clinical diagnosis based on characteristic migratory, erythematous, well-demarcated patches with raised white borders
  \item Inspection of the entire oral mucosa for associated lesions (psoriasis, lichen planus)
  \item Biopsy reserved for atypical presentations (non-migratory, painful, or persistent): shows spongiform pustules of Kogoj (similar to pustular psoriasis)
  \item KOH preparation to exclude candidiasis (if white border is scrapable)
  \item Patch testing if contact allergy suspected (rare)
\end{itemize}

\section*{Treatment Options}
\begin{itemize}
  \item No treatment for asymptomatic cases -- reassurance is the main intervention
  \item Topical corticosteroids (0.1\% triamcinolone dental paste or clobetasol gel) for painful exacerbations
  \item Topical calcineurin inhibitors (tacrolimus 0.1\% ointment) for refractory symptoms
  \item Topical antihistamine mouth rinses (diphenhydramine elixir swish and spit)
  \item Anesthetic rinses for pre-meal pain relief
  \item Zinc supplementation (220 mg zinc sulfate daily) has limited evidence of benefit
  \item Avoidance of triggers identified by you
\end{itemize}

\section*{References}
\begin{thebibliography}{99}
\bibitem{ref1} Reamy BV, Derby R, Bunt CW. "Common Tongue Conditions in Primary Care." Am Fam Physician. 2010;81(5):627--634.
\bibitem{ref2} Assimakopoulos D, Patrikakos G, Fotika C, et al. "Benign Migratory Glossitis or Geographic Tongue: An Enigmatic Oral Lesion." Am J Med. 2002;113(9):751--755.
\bibitem{ref3} Zadik Y, Drucker S, Pallmon S. "Migratory Glossitis (Geographic Tongue): A Review." Quintessence Int. 2011;42(8):687--692.
\bibitem{ref4} Neville BW, Damm DD, Allen CM, et al. "Oral and Maxillofacial Pathology." 4th ed. Elsevier; 2016.
\bibitem{ref5} Scully C. "Oral and Maxillofacial Medicine." 3rd ed. Churchill Livingstone, 2013.
\end{thebibliography}
